% =============================================================================
%  Phase 2: Data Collection & Preparation (preprocessing documentation)
%  Limit: 1-2 pages or a well-structured README. Rubric: docs/spec/phase2_details.pdf
%  Source text: reports/phase2_report.md. Numbers verified against
%  notebooks/data_preparation.ipynb and the Parquet outputs.
%  Lines marked "% CHECK:" are review notes for the authors, not report text.
%  Build: latexmk -pdf -outdir=build main.tex
% =============================================================================
\documentclass[11pt,letterpaper]{article}
% =============================================================================
%  preamble.tex -- shared styling for all course report deliverables
%  Usage (from reports/phaseN/main.tex):
%      \documentclass[11pt,letterpaper]{article}
%      % =============================================================================
%  preamble.tex -- shared styling for all course report deliverables
%  Usage (from reports/phaseN/main.tex):
%      \documentclass[11pt,letterpaper]{article}
%      % =============================================================================
%  preamble.tex -- shared styling for all course report deliverables
%  Usage (from reports/phaseN/main.tex):
%      \documentclass[11pt,letterpaper]{article}
%      \input{../_template/preamble}
%  Engine: pdfLaTeX + biber (latexmk handles the run order).
%  Do not put content here. Do not use inline colour values in documents;
%  use the named palette below.
% =============================================================================

\makeatletter   % internal \@ names are used throughout; closed at end of file


% -----------------------------------------------------------------------------
%  1. PACKAGES: encoding and language
% -----------------------------------------------------------------------------
\usepackage[T1]{fontenc}          % proper hyphenation of accented glyphs, copyable PDF text
\usepackage[utf8]{inputenc}       % default on modern kernels; kept for older installs
\usepackage[english]{babel}       % hyphenation patterns
\usepackage{csquotes}             % language-aware quotes; biblatex expects it with babel


% -----------------------------------------------------------------------------
%  2. COLOURS: the only place colour values are defined
%     Colour is used for headings, rules, table header rows and links. Nothing else.
% -----------------------------------------------------------------------------
\usepackage[table]{xcolor}        % 'table' loads colortbl for \rowcolor

\definecolor{brandNavy}{HTML}{1F3A5F}    % primary: headings, title, rules
\definecolor{brandAccent}{HTML}{23707D}  % accent: links, title kicker (muted teal)
\definecolor{inkText}{HTML}{1A1D21}      % body text (near-black, softer than pure black)
\definecolor{mutedText}{HTML}{5B6570}    % header/footer, metadata line
\definecolor{ruleGray}{HTML}{B8BEC6}     % light rules (running header)
\definecolor{tableHeadBg}{HTML}{E6ECF2}  % table header row tint (navy at ~10%)


% -----------------------------------------------------------------------------
%  3. FONTS
%     Body:     XCharter, a Type 1 extension of Bitstream Charter. Sturdy, open letterforms
%               that stay legible at 11pt on screen and in print, and a moderate set width.
%     Math:     newtxmath with the xcharter option, so math matches the body.
%     Headings: Source Sans Pro, a humanist sans that pairs with Charter.
%     Mono:     Inconsolata (zi4) for file names and code, scaled to the body x-height.
% -----------------------------------------------------------------------------
\usepackage[osf=false]{XCharter}
\usepackage[xcharter,vvarbb]{newtxmath}
\IfFileExists{sourcesans.sty}            % package renamed in 2022; fall back on older installs
  {\usepackage[semibold]{sourcesans}}    % sets \sffamily only; body stays serif
  {\usepackage[semibold]{sourcesanspro}}
\usepackage[scaled=0.95,varqu]{zi4}      % sets \ttfamily only

\usepackage[final]{microtype}            % protrusion + font expansion: fewer bad breaks, denser text


% -----------------------------------------------------------------------------
%  4. LAYOUT
% -----------------------------------------------------------------------------
\usepackage[letterpaper,margin=1in,headheight=14pt,headsep=12pt,footskip=28pt]{geometry}

\color{inkText}                          % default text colour
\setlength{\parindent}{1.2em}            % indented paragraphs, no gap: densest readable setting
\setlength{\parskip}{0pt}
\linespread{1.04}                        % slight extra leading; Charter is dark and wide

\clubpenalty=10000                       % no widows or orphans
\widowpenalty=10000
\displaywidowpenalty=10000

\usepackage{enumitem}                    % tight lists
\setlist{itemsep=1pt,parsep=0pt,topsep=3pt,partopsep=0pt}
\setlist[itemize]{leftmargin=1.4em}
\setlist[enumerate]{leftmargin=1.8em}
\setlist[description]{leftmargin=1.2em,font=\sffamily\bfseries\color{brandNavy}}

\usepackage{graphicx}
\graphicspath{{../figures/}}             % figures come from reports/figures/ (built by src/.../viz)

\usepackage[font=small,labelfont=bf,skip=4pt]{caption}
\captionsetup[table]{position=top}       % table captions go above the table

\setlength{\textfloatsep}{10pt plus 2pt minus 2pt}   % tighter space around floats
\setlength{\floatsep}{8pt plus 2pt minus 2pt}
\setlength{\intextsep}{8pt plus 2pt minus 2pt}


% -----------------------------------------------------------------------------
%  5. HEADINGS
%     Section:       sans semibold, large, navy, numbered
%     Subsection:    sans semibold, body size, navy
%     Subsubsection: sans semibold, body size, near-black
%     Paragraph:     run-in sans semibold
% -----------------------------------------------------------------------------
\usepackage[explicit]{titlesec}

\titleformat{\section}
  {\sffamily\large\bfseries\color{brandNavy}}{\thesection}{0.6em}{#1}
\titleformat{\subsection}
  {\sffamily\normalsize\bfseries\color{brandNavy}}{\thesubsection}{0.6em}{#1}
\titleformat{\subsubsection}
  {\sffamily\normalsize\bfseries\color{inkText}}{\thesubsubsection}{0.6em}{#1}
\titleformat{\paragraph}[runin]
  {\sffamily\normalsize\bfseries\color{inkText}}{}{0pt}{#1.}

%                         left  before                 after
\titlespacing*{\section}      {0pt}{1.6ex plus .4ex minus .2ex}{0.6ex plus .1ex}
\titlespacing*{\subsection}   {0pt}{1.2ex plus .3ex minus .2ex}{0.4ex plus .1ex}
\titlespacing*{\subsubsection}{0pt}{1.0ex plus .3ex minus .2ex}{0.3ex}
\titlespacing*{\paragraph}    {0pt}{0.8ex plus .2ex minus .1ex}{0.6em}


% -----------------------------------------------------------------------------
%  6. TABLES
%     booktabs rules only (no vertical rules), tinted header row, tabularx for wrapping.
%     Pattern:
%        \tableheadrow \colhead{A} & \colhead{B} \\   <- replaces \toprule + header row
%        \tableheadend                                 <- replaces \midrule under the header
%        body rows ...
%        \bottomrule
%     The header rules have zero separation so the tint meets them with no white gap;
%     \colhead adds symmetric padding instead. Body rows keep standard booktabs spacing.
% -----------------------------------------------------------------------------
\usepackage{array}
\usepackage{booktabs}
\usepackage{tabularx}
\usepackage{ragged2e}

\renewcommand{\arraystretch}{1.15}       % a little air between wrapped body rows
\arrayrulecolor{brandNavy}
\newlength{\cellpad}\setlength{\cellpad}{2.5pt}   % header row padding above and below

% L: ragged-right X column (allows hyphenation, avoids stretched word spacing in narrow cells)
% Y{f}: same, with relative width f. In one table the f values must sum to the number of X/Y columns.
\newcolumntype{L}{>{\RaggedRight\arraybackslash}X}
\newcolumntype{Y}[1]{>{\hsize=#1\hsize\RaggedRight\arraybackslash}X}

\newcommand{\tableheadrow}{\specialrule{\heavyrulewidth}{\abovetopsep}{0pt}\rowcolor{tableHeadBg}}
\newcommand{\tableheadend}{\specialrule{\lightrulewidth}{0pt}{\belowrulesep}}
\newcommand{\colhead}[1]{%
  \rule{0pt}{\dimexpr0.7\baselineskip+\cellpad\relax}%               top padding strut
  {\sffamily\bfseries\color{brandNavy}#1}%
  \rule[\dimexpr-0.3\baselineskip-\cellpad\relax]{0pt}{\dimexpr0.3\baselineskip+\cellpad\relax}}% bottom


% -----------------------------------------------------------------------------
%  7. HEADER / FOOTER
%     Header: short project title (left), phase label (right), thin gray rule.
%     Footer: page number (centre). First page: footer only (title block replaces header).
% -----------------------------------------------------------------------------
\usepackage{fancyhdr}
\pagestyle{fancy}
\fancyhf{}
\fancyhead[L]{\sffamily\small\color{mutedText}\@shorttitle}
\fancyhead[R]{\sffamily\small\color{mutedText}\@phaselabel}
\fancyfoot[C]{\sffamily\small\color{mutedText}\thepage}
\renewcommand{\headrulewidth}{0.4pt}
\renewcommand{\headrule}{{\color{ruleGray}%
  \hrule height\headrulewidth width\headwidth \vskip-\headrulewidth}}

\fancypagestyle{plain}{%
  \fancyhf{}%
  \fancyfoot[C]{\sffamily\small\color{mutedText}\thepage}%
  \renewcommand{\headrulewidth}{0pt}%
}


% -----------------------------------------------------------------------------
%  8. BIBLIOGRAPHY: biblatex + biber, numeric-comp
%     Numeric keeps in-text citations short (matters under a page limit); '-comp'
%     compresses [1,2,3] to [1-3]. Ordered by first citation.
% -----------------------------------------------------------------------------
\usepackage[backend=biber,style=numeric-comp,sorting=none,maxbibnames=6,
            giveninits=true,url=true,doi=true,isbn=false,
            date=iso,urldate=iso,seconds=true]{biblatex}   % ISO dates: unambiguous across locales
\renewcommand*{\bibfont}{\small}
\setlength{\bibitemsep}{2pt}


% -----------------------------------------------------------------------------
%  9. METADATA MACROS
%     Set in each document before \begin{document}:
%       \projecttitle[Short header title]{Full title}
%       \phaselabel{...}  \authors{...}  \course{...}  \reportdate{...}
% -----------------------------------------------------------------------------
\newcommand{\@projecttitle}{[PROJECT TITLE]}
\newcommand{\@shorttitle}{[SHORT TITLE]}
\newcommand{\@phaselabel}{[PHASE LABEL]}
\newcommand{\@authors}{[AUTHORS]}
\newcommand{\@course}{[COURSE]}
\newcommand{\@reportdate}{[DATE]}

\newcommand{\projecttitle}[2][]{%
  \renewcommand{\@projecttitle}{#2}%
  \if\relax\detokenize{#1}\relax\renewcommand{\@shorttitle}{#2}%
  \else\renewcommand{\@shorttitle}{#1}\fi}
\newcommand{\phaselabel}[1]{\renewcommand{\@phaselabel}{#1}}
\newcommand{\authors}[1]{\renewcommand{\@authors}{#1}}
\newcommand{\course}[1]{\renewcommand{\@course}{#1}}
\newcommand{\reportdate}[1]{\renewcommand{\@reportdate}{#1}}

% Compact title block (not a title page). Call right after \begin{document}.
\newcommand{\maketitleblock}{%
  \thispagestyle{plain}%
  \begingroup\raggedright
  {\sffamily\small\bfseries\color{brandAccent}\textls[60]{\MakeUppercase{\@phaselabel}}\par}
  \vspace{3pt}
  {\sffamily\LARGE\bfseries\color{brandNavy}\@projecttitle\par}
  \vspace{6pt}
  {\normalsize\@authors\par}
  \vspace{2pt}
  {\small\color{mutedText}\@course\ \textbar\ \@reportdate\par}
  \vspace{2pt}
  {\color{brandNavy}\rule{\linewidth}{0.8pt}\par}
  \endgroup
  \vspace{4pt}%
}

% Inline code / file names
\newcommand{\code}[1]{\texttt{#1}}


% -----------------------------------------------------------------------------
%  10. HYPERLINKS (load last)
% -----------------------------------------------------------------------------
\usepackage[colorlinks=true,linkcolor=brandNavy,citecolor=brandAccent,
            urlcolor=brandAccent,filecolor=brandAccent,bookmarksnumbered=true,
            pdfstartview=FitH]{hyperref}
\urlstyle{same}                          % URLs in body font: narrower than monospace

\AtBeginDocument{\hypersetup{pdftitle={\@projecttitle},pdfauthor={\@authors},
                             pdfsubject={\@phaselabel}}}
\makeatother

%  Engine: pdfLaTeX + biber (latexmk handles the run order).
%  Do not put content here. Do not use inline colour values in documents;
%  use the named palette below.
% =============================================================================

\makeatletter   % internal \@ names are used throughout; closed at end of file


% -----------------------------------------------------------------------------
%  1. PACKAGES: encoding and language
% -----------------------------------------------------------------------------
\usepackage[T1]{fontenc}          % proper hyphenation of accented glyphs, copyable PDF text
\usepackage[utf8]{inputenc}       % default on modern kernels; kept for older installs
\usepackage[english]{babel}       % hyphenation patterns
\usepackage{csquotes}             % language-aware quotes; biblatex expects it with babel


% -----------------------------------------------------------------------------
%  2. COLOURS: the only place colour values are defined
%     Colour is used for headings, rules, table header rows and links. Nothing else.
% -----------------------------------------------------------------------------
\usepackage[table]{xcolor}        % 'table' loads colortbl for \rowcolor

\definecolor{brandNavy}{HTML}{1F3A5F}    % primary: headings, title, rules
\definecolor{brandAccent}{HTML}{23707D}  % accent: links, title kicker (muted teal)
\definecolor{inkText}{HTML}{1A1D21}      % body text (near-black, softer than pure black)
\definecolor{mutedText}{HTML}{5B6570}    % header/footer, metadata line
\definecolor{ruleGray}{HTML}{B8BEC6}     % light rules (running header)
\definecolor{tableHeadBg}{HTML}{E6ECF2}  % table header row tint (navy at ~10%)


% -----------------------------------------------------------------------------
%  3. FONTS
%     Body:     XCharter, a Type 1 extension of Bitstream Charter. Sturdy, open letterforms
%               that stay legible at 11pt on screen and in print, and a moderate set width.
%     Math:     newtxmath with the xcharter option, so math matches the body.
%     Headings: Source Sans Pro, a humanist sans that pairs with Charter.
%     Mono:     Inconsolata (zi4) for file names and code, scaled to the body x-height.
% -----------------------------------------------------------------------------
\usepackage[osf=false]{XCharter}
\usepackage[xcharter,vvarbb]{newtxmath}
\IfFileExists{sourcesans.sty}            % package renamed in 2022; fall back on older installs
  {\usepackage[semibold]{sourcesans}}    % sets \sffamily only; body stays serif
  {\usepackage[semibold]{sourcesanspro}}
\usepackage[scaled=0.95,varqu]{zi4}      % sets \ttfamily only

\usepackage[final]{microtype}            % protrusion + font expansion: fewer bad breaks, denser text


% -----------------------------------------------------------------------------
%  4. LAYOUT
% -----------------------------------------------------------------------------
\usepackage[letterpaper,margin=1in,headheight=14pt,headsep=12pt,footskip=28pt]{geometry}

\color{inkText}                          % default text colour
\setlength{\parindent}{1.2em}            % indented paragraphs, no gap: densest readable setting
\setlength{\parskip}{0pt}
\linespread{1.04}                        % slight extra leading; Charter is dark and wide

\clubpenalty=10000                       % no widows or orphans
\widowpenalty=10000
\displaywidowpenalty=10000

\usepackage{enumitem}                    % tight lists
\setlist{itemsep=1pt,parsep=0pt,topsep=3pt,partopsep=0pt}
\setlist[itemize]{leftmargin=1.4em}
\setlist[enumerate]{leftmargin=1.8em}
\setlist[description]{leftmargin=1.2em,font=\sffamily\bfseries\color{brandNavy}}

\usepackage{graphicx}
\graphicspath{{../figures/}}             % figures come from reports/figures/ (built by src/.../viz)

\usepackage[font=small,labelfont=bf,skip=4pt]{caption}
\captionsetup[table]{position=top}       % table captions go above the table

\setlength{\textfloatsep}{10pt plus 2pt minus 2pt}   % tighter space around floats
\setlength{\floatsep}{8pt plus 2pt minus 2pt}
\setlength{\intextsep}{8pt plus 2pt minus 2pt}


% -----------------------------------------------------------------------------
%  5. HEADINGS
%     Section:       sans semibold, large, navy, numbered
%     Subsection:    sans semibold, body size, navy
%     Subsubsection: sans semibold, body size, near-black
%     Paragraph:     run-in sans semibold
% -----------------------------------------------------------------------------
\usepackage[explicit]{titlesec}

\titleformat{\section}
  {\sffamily\large\bfseries\color{brandNavy}}{\thesection}{0.6em}{#1}
\titleformat{\subsection}
  {\sffamily\normalsize\bfseries\color{brandNavy}}{\thesubsection}{0.6em}{#1}
\titleformat{\subsubsection}
  {\sffamily\normalsize\bfseries\color{inkText}}{\thesubsubsection}{0.6em}{#1}
\titleformat{\paragraph}[runin]
  {\sffamily\normalsize\bfseries\color{inkText}}{}{0pt}{#1.}

%                         left  before                 after
\titlespacing*{\section}      {0pt}{1.6ex plus .4ex minus .2ex}{0.6ex plus .1ex}
\titlespacing*{\subsection}   {0pt}{1.2ex plus .3ex minus .2ex}{0.4ex plus .1ex}
\titlespacing*{\subsubsection}{0pt}{1.0ex plus .3ex minus .2ex}{0.3ex}
\titlespacing*{\paragraph}    {0pt}{0.8ex plus .2ex minus .1ex}{0.6em}


% -----------------------------------------------------------------------------
%  6. TABLES
%     booktabs rules only (no vertical rules), tinted header row, tabularx for wrapping.
%     Pattern:
%        \tableheadrow \colhead{A} & \colhead{B} \\   <- replaces \toprule + header row
%        \tableheadend                                 <- replaces \midrule under the header
%        body rows ...
%        \bottomrule
%     The header rules have zero separation so the tint meets them with no white gap;
%     \colhead adds symmetric padding instead. Body rows keep standard booktabs spacing.
% -----------------------------------------------------------------------------
\usepackage{array}
\usepackage{booktabs}
\usepackage{tabularx}
\usepackage{ragged2e}

\renewcommand{\arraystretch}{1.15}       % a little air between wrapped body rows
\arrayrulecolor{brandNavy}
\newlength{\cellpad}\setlength{\cellpad}{2.5pt}   % header row padding above and below

% L: ragged-right X column (allows hyphenation, avoids stretched word spacing in narrow cells)
% Y{f}: same, with relative width f. In one table the f values must sum to the number of X/Y columns.
\newcolumntype{L}{>{\RaggedRight\arraybackslash}X}
\newcolumntype{Y}[1]{>{\hsize=#1\hsize\RaggedRight\arraybackslash}X}

\newcommand{\tableheadrow}{\specialrule{\heavyrulewidth}{\abovetopsep}{0pt}\rowcolor{tableHeadBg}}
\newcommand{\tableheadend}{\specialrule{\lightrulewidth}{0pt}{\belowrulesep}}
\newcommand{\colhead}[1]{%
  \rule{0pt}{\dimexpr0.7\baselineskip+\cellpad\relax}%               top padding strut
  {\sffamily\bfseries\color{brandNavy}#1}%
  \rule[\dimexpr-0.3\baselineskip-\cellpad\relax]{0pt}{\dimexpr0.3\baselineskip+\cellpad\relax}}% bottom


% -----------------------------------------------------------------------------
%  7. HEADER / FOOTER
%     Header: short project title (left), phase label (right), thin gray rule.
%     Footer: page number (centre). First page: footer only (title block replaces header).
% -----------------------------------------------------------------------------
\usepackage{fancyhdr}
\pagestyle{fancy}
\fancyhf{}
\fancyhead[L]{\sffamily\small\color{mutedText}\@shorttitle}
\fancyhead[R]{\sffamily\small\color{mutedText}\@phaselabel}
\fancyfoot[C]{\sffamily\small\color{mutedText}\thepage}
\renewcommand{\headrulewidth}{0.4pt}
\renewcommand{\headrule}{{\color{ruleGray}%
  \hrule height\headrulewidth width\headwidth \vskip-\headrulewidth}}

\fancypagestyle{plain}{%
  \fancyhf{}%
  \fancyfoot[C]{\sffamily\small\color{mutedText}\thepage}%
  \renewcommand{\headrulewidth}{0pt}%
}


% -----------------------------------------------------------------------------
%  8. BIBLIOGRAPHY: biblatex + biber, numeric-comp
%     Numeric keeps in-text citations short (matters under a page limit); '-comp'
%     compresses [1,2,3] to [1-3]. Ordered by first citation.
% -----------------------------------------------------------------------------
\usepackage[backend=biber,style=numeric-comp,sorting=none,maxbibnames=6,
            giveninits=true,url=true,doi=true,isbn=false,
            date=iso,urldate=iso,seconds=true]{biblatex}   % ISO dates: unambiguous across locales
\renewcommand*{\bibfont}{\small}
\setlength{\bibitemsep}{2pt}


% -----------------------------------------------------------------------------
%  9. METADATA MACROS
%     Set in each document before \begin{document}:
%       \projecttitle[Short header title]{Full title}
%       \phaselabel{...}  \authors{...}  \course{...}  \reportdate{...}
% -----------------------------------------------------------------------------
\newcommand{\@projecttitle}{[PROJECT TITLE]}
\newcommand{\@shorttitle}{[SHORT TITLE]}
\newcommand{\@phaselabel}{[PHASE LABEL]}
\newcommand{\@authors}{[AUTHORS]}
\newcommand{\@course}{[COURSE]}
\newcommand{\@reportdate}{[DATE]}

\newcommand{\projecttitle}[2][]{%
  \renewcommand{\@projecttitle}{#2}%
  \if\relax\detokenize{#1}\relax\renewcommand{\@shorttitle}{#2}%
  \else\renewcommand{\@shorttitle}{#1}\fi}
\newcommand{\phaselabel}[1]{\renewcommand{\@phaselabel}{#1}}
\newcommand{\authors}[1]{\renewcommand{\@authors}{#1}}
\newcommand{\course}[1]{\renewcommand{\@course}{#1}}
\newcommand{\reportdate}[1]{\renewcommand{\@reportdate}{#1}}

% Compact title block (not a title page). Call right after \begin{document}.
\newcommand{\maketitleblock}{%
  \thispagestyle{plain}%
  \begingroup\raggedright
  {\sffamily\small\bfseries\color{brandAccent}\textls[60]{\MakeUppercase{\@phaselabel}}\par}
  \vspace{3pt}
  {\sffamily\LARGE\bfseries\color{brandNavy}\@projecttitle\par}
  \vspace{6pt}
  {\normalsize\@authors\par}
  \vspace{2pt}
  {\small\color{mutedText}\@course\ \textbar\ \@reportdate\par}
  \vspace{2pt}
  {\color{brandNavy}\rule{\linewidth}{0.8pt}\par}
  \endgroup
  \vspace{4pt}%
}

% Inline code / file names
\newcommand{\code}[1]{\texttt{#1}}


% -----------------------------------------------------------------------------
%  10. HYPERLINKS (load last)
% -----------------------------------------------------------------------------
\usepackage[colorlinks=true,linkcolor=brandNavy,citecolor=brandAccent,
            urlcolor=brandAccent,filecolor=brandAccent,bookmarksnumbered=true,
            pdfstartview=FitH]{hyperref}
\urlstyle{same}                          % URLs in body font: narrower than monospace

\AtBeginDocument{\hypersetup{pdftitle={\@projecttitle},pdfauthor={\@authors},
                             pdfsubject={\@phaselabel}}}
\makeatother

%  Engine: pdfLaTeX + biber (latexmk handles the run order).
%  Do not put content here. Do not use inline colour values in documents;
%  use the named palette below.
% =============================================================================

\makeatletter   % internal \@ names are used throughout; closed at end of file


% -----------------------------------------------------------------------------
%  1. PACKAGES: encoding and language
% -----------------------------------------------------------------------------
\usepackage[T1]{fontenc}          % proper hyphenation of accented glyphs, copyable PDF text
\usepackage[utf8]{inputenc}       % default on modern kernels; kept for older installs
\usepackage[english]{babel}       % hyphenation patterns
\usepackage{csquotes}             % language-aware quotes; biblatex expects it with babel


% -----------------------------------------------------------------------------
%  2. COLOURS: the only place colour values are defined
%     Colour is used for headings, rules, table header rows and links. Nothing else.
% -----------------------------------------------------------------------------
\usepackage[table]{xcolor}        % 'table' loads colortbl for \rowcolor

\definecolor{brandNavy}{HTML}{1F3A5F}    % primary: headings, title, rules
\definecolor{brandAccent}{HTML}{23707D}  % accent: links, title kicker (muted teal)
\definecolor{inkText}{HTML}{1A1D21}      % body text (near-black, softer than pure black)
\definecolor{mutedText}{HTML}{5B6570}    % header/footer, metadata line
\definecolor{ruleGray}{HTML}{B8BEC6}     % light rules (running header)
\definecolor{tableHeadBg}{HTML}{E6ECF2}  % table header row tint (navy at ~10%)


% -----------------------------------------------------------------------------
%  3. FONTS
%     Body:     XCharter, a Type 1 extension of Bitstream Charter. Sturdy, open letterforms
%               that stay legible at 11pt on screen and in print, and a moderate set width.
%     Math:     newtxmath with the xcharter option, so math matches the body.
%     Headings: Source Sans Pro, a humanist sans that pairs with Charter.
%     Mono:     Inconsolata (zi4) for file names and code, scaled to the body x-height.
% -----------------------------------------------------------------------------
\usepackage[osf=false]{XCharter}
\usepackage[xcharter,vvarbb]{newtxmath}
\IfFileExists{sourcesans.sty}            % package renamed in 2022; fall back on older installs
  {\usepackage[semibold]{sourcesans}}    % sets \sffamily only; body stays serif
  {\usepackage[semibold]{sourcesanspro}}
\usepackage[scaled=0.95,varqu]{zi4}      % sets \ttfamily only

\usepackage[final]{microtype}            % protrusion + font expansion: fewer bad breaks, denser text


% -----------------------------------------------------------------------------
%  4. LAYOUT
% -----------------------------------------------------------------------------
\usepackage[letterpaper,margin=1in,headheight=14pt,headsep=12pt,footskip=28pt]{geometry}

\color{inkText}                          % default text colour
\setlength{\parindent}{1.2em}            % indented paragraphs, no gap: densest readable setting
\setlength{\parskip}{0pt}
\linespread{1.04}                        % slight extra leading; Charter is dark and wide

\clubpenalty=10000                       % no widows or orphans
\widowpenalty=10000
\displaywidowpenalty=10000

\usepackage{enumitem}                    % tight lists
\setlist{itemsep=1pt,parsep=0pt,topsep=3pt,partopsep=0pt}
\setlist[itemize]{leftmargin=1.4em}
\setlist[enumerate]{leftmargin=1.8em}
\setlist[description]{leftmargin=1.2em,font=\sffamily\bfseries\color{brandNavy}}

\usepackage{graphicx}
\graphicspath{{../figures/}}             % figures come from reports/figures/ (built by src/.../viz)

\usepackage[font=small,labelfont=bf,skip=4pt]{caption}
\captionsetup[table]{position=top}       % table captions go above the table

\setlength{\textfloatsep}{10pt plus 2pt minus 2pt}   % tighter space around floats
\setlength{\floatsep}{8pt plus 2pt minus 2pt}
\setlength{\intextsep}{8pt plus 2pt minus 2pt}


% -----------------------------------------------------------------------------
%  5. HEADINGS
%     Section:       sans semibold, large, navy, numbered
%     Subsection:    sans semibold, body size, navy
%     Subsubsection: sans semibold, body size, near-black
%     Paragraph:     run-in sans semibold
% -----------------------------------------------------------------------------
\usepackage[explicit]{titlesec}

\titleformat{\section}
  {\sffamily\large\bfseries\color{brandNavy}}{\thesection}{0.6em}{#1}
\titleformat{\subsection}
  {\sffamily\normalsize\bfseries\color{brandNavy}}{\thesubsection}{0.6em}{#1}
\titleformat{\subsubsection}
  {\sffamily\normalsize\bfseries\color{inkText}}{\thesubsubsection}{0.6em}{#1}
\titleformat{\paragraph}[runin]
  {\sffamily\normalsize\bfseries\color{inkText}}{}{0pt}{#1.}

%                         left  before                 after
\titlespacing*{\section}      {0pt}{1.6ex plus .4ex minus .2ex}{0.6ex plus .1ex}
\titlespacing*{\subsection}   {0pt}{1.2ex plus .3ex minus .2ex}{0.4ex plus .1ex}
\titlespacing*{\subsubsection}{0pt}{1.0ex plus .3ex minus .2ex}{0.3ex}
\titlespacing*{\paragraph}    {0pt}{0.8ex plus .2ex minus .1ex}{0.6em}


% -----------------------------------------------------------------------------
%  6. TABLES
%     booktabs rules only (no vertical rules), tinted header row, tabularx for wrapping.
%     Pattern:
%        \tableheadrow \colhead{A} & \colhead{B} \\   <- replaces \toprule + header row
%        \tableheadend                                 <- replaces \midrule under the header
%        body rows ...
%        \bottomrule
%     The header rules have zero separation so the tint meets them with no white gap;
%     \colhead adds symmetric padding instead. Body rows keep standard booktabs spacing.
% -----------------------------------------------------------------------------
\usepackage{array}
\usepackage{booktabs}
\usepackage{tabularx}
\usepackage{ragged2e}

\renewcommand{\arraystretch}{1.15}       % a little air between wrapped body rows
\arrayrulecolor{brandNavy}
\newlength{\cellpad}\setlength{\cellpad}{2.5pt}   % header row padding above and below

% L: ragged-right X column (allows hyphenation, avoids stretched word spacing in narrow cells)
% Y{f}: same, with relative width f. In one table the f values must sum to the number of X/Y columns.
\newcolumntype{L}{>{\RaggedRight\arraybackslash}X}
\newcolumntype{Y}[1]{>{\hsize=#1\hsize\RaggedRight\arraybackslash}X}

\newcommand{\tableheadrow}{\specialrule{\heavyrulewidth}{\abovetopsep}{0pt}\rowcolor{tableHeadBg}}
\newcommand{\tableheadend}{\specialrule{\lightrulewidth}{0pt}{\belowrulesep}}
\newcommand{\colhead}[1]{%
  \rule{0pt}{\dimexpr0.7\baselineskip+\cellpad\relax}%               top padding strut
  {\sffamily\bfseries\color{brandNavy}#1}%
  \rule[\dimexpr-0.3\baselineskip-\cellpad\relax]{0pt}{\dimexpr0.3\baselineskip+\cellpad\relax}}% bottom


% -----------------------------------------------------------------------------
%  7. HEADER / FOOTER
%     Header: short project title (left), phase label (right), thin gray rule.
%     Footer: page number (centre). First page: footer only (title block replaces header).
% -----------------------------------------------------------------------------
\usepackage{fancyhdr}
\pagestyle{fancy}
\fancyhf{}
\fancyhead[L]{\sffamily\small\color{mutedText}\@shorttitle}
\fancyhead[R]{\sffamily\small\color{mutedText}\@phaselabel}
\fancyfoot[C]{\sffamily\small\color{mutedText}\thepage}
\renewcommand{\headrulewidth}{0.4pt}
\renewcommand{\headrule}{{\color{ruleGray}%
  \hrule height\headrulewidth width\headwidth \vskip-\headrulewidth}}

\fancypagestyle{plain}{%
  \fancyhf{}%
  \fancyfoot[C]{\sffamily\small\color{mutedText}\thepage}%
  \renewcommand{\headrulewidth}{0pt}%
}


% -----------------------------------------------------------------------------
%  8. BIBLIOGRAPHY: biblatex + biber, numeric-comp
%     Numeric keeps in-text citations short (matters under a page limit); '-comp'
%     compresses [1,2,3] to [1-3]. Ordered by first citation.
% -----------------------------------------------------------------------------
\usepackage[backend=biber,style=numeric-comp,sorting=none,maxbibnames=6,
            giveninits=true,url=true,doi=true,isbn=false,
            date=iso,urldate=iso,seconds=true]{biblatex}   % ISO dates: unambiguous across locales
\renewcommand*{\bibfont}{\small}
\setlength{\bibitemsep}{2pt}


% -----------------------------------------------------------------------------
%  9. METADATA MACROS
%     Set in each document before \begin{document}:
%       \projecttitle[Short header title]{Full title}
%       \phaselabel{...}  \authors{...}  \course{...}  \reportdate{...}
% -----------------------------------------------------------------------------
\newcommand{\@projecttitle}{[PROJECT TITLE]}
\newcommand{\@shorttitle}{[SHORT TITLE]}
\newcommand{\@phaselabel}{[PHASE LABEL]}
\newcommand{\@authors}{[AUTHORS]}
\newcommand{\@course}{[COURSE]}
\newcommand{\@reportdate}{[DATE]}

\newcommand{\projecttitle}[2][]{%
  \renewcommand{\@projecttitle}{#2}%
  \if\relax\detokenize{#1}\relax\renewcommand{\@shorttitle}{#2}%
  \else\renewcommand{\@shorttitle}{#1}\fi}
\newcommand{\phaselabel}[1]{\renewcommand{\@phaselabel}{#1}}
\newcommand{\authors}[1]{\renewcommand{\@authors}{#1}}
\newcommand{\course}[1]{\renewcommand{\@course}{#1}}
\newcommand{\reportdate}[1]{\renewcommand{\@reportdate}{#1}}

% Compact title block (not a title page). Call right after \begin{document}.
\newcommand{\maketitleblock}{%
  \thispagestyle{plain}%
  \begingroup\raggedright
  {\sffamily\small\bfseries\color{brandAccent}\textls[60]{\MakeUppercase{\@phaselabel}}\par}
  \vspace{3pt}
  {\sffamily\LARGE\bfseries\color{brandNavy}\@projecttitle\par}
  \vspace{6pt}
  {\normalsize\@authors\par}
  \vspace{2pt}
  {\small\color{mutedText}\@course\ \textbar\ \@reportdate\par}
  \vspace{2pt}
  {\color{brandNavy}\rule{\linewidth}{0.8pt}\par}
  \endgroup
  \vspace{4pt}%
}

% Inline code / file names
\newcommand{\code}[1]{\texttt{#1}}


% -----------------------------------------------------------------------------
%  10. HYPERLINKS (load last)
% -----------------------------------------------------------------------------
\usepackage[colorlinks=true,linkcolor=brandNavy,citecolor=brandAccent,
            urlcolor=brandAccent,filecolor=brandAccent,bookmarksnumbered=true,
            pdfstartview=FitH]{hyperref}
\urlstyle{same}                          % URLs in body font: narrower than monospace

\AtBeginDocument{\hypersetup{pdftitle={\@projecttitle},pdfauthor={\@authors},
                             pdfsubject={\@phaselabel}}}
\makeatother

\addbibresource{references.bib}

% ---- Metadata ---------------------------------------------------------------
\projecttitle[Bill Forecasting]{Household Electricity Consumption and Bill Forecasting}
\phaselabel{Phase 2 \textperiodcentered{} Data Collection \& Preparation}
\authors{Akinbobola Akin (3784664), Adejuwon Oluwasegun Isaac (3806766), Edehi Andrew Igue (3778767)}
\course{TME6014 -- Data Analytics}
\reportdate{\today}

\begin{document}
\maketitleblock

% =============================================================================
\section{Data Collection}
\label{sec:collection}

The project uses the \emph{Smart Meters in London} dataset, originating from
the Low Carbon London project. The dataset contains half-hourly
electricity consumption readings collected between November 2011 and
February 2014. It was downloaded from Kaggle~\cite{kaggle_smartmeters}
using Python's kagglehub library (version 21 of the dataset). The
downloaded files include 112 electricity consumption CSV files,
household information, weather observations, and UK bank holiday
records. Electricity tariff information was obtained from the London
Datastore~\cite{datastore_smartmeters} through Tariffs.xlsx, included in the project repository.
% CHECK: rubric asks for each source's URL, access method, description and
% retrieval date. Retrieval date is only in the bib (2026-11-01, from file
% timestamps; confirm). Weather (Dark Sky), household and holiday files have
% no description of their own.

The original dataset contained \textbf{167,817,021 readings from 5,566
households}. DuckDB was used to process the large dataset efficiently,
while CSV files were converted to Parquet format for improved storage
and query performance.

% =============================================================================
\section{Data Preprocessing}
\label{sec:preprocessing}

\subsection{Data Acquisition Validation}
\label{sec:validation}

The downloaded CSV files were examined for readability, column
consistency, and record integrity. After conversion to Parquet, row
counts were compared to confirm that all 167,817,021 records were
preserved. Initial profiling also examined household identifiers,
timestamps, and electricity consumption values.

\subsection{Missing Data Handling}
\label{sec:missing}

Initial validation identified 5,560 records containing invalid Null
energy readings, which were removed. Missing observations were
subsequently detected by comparing available readings against the
expected 30-minute intervals, revealing 456,658 missing intervals
(0.271\%). Extended zero-consumption periods lasting at least 24 hours
were also treated as missing.

A \textbf{90\% monthly coverage threshold} was adopted to exclude
insufficiently complete household-months. Short gaps of up to two hours
were filled using linear interpolation, while remaining gaps were
estimated using household-specific average consumption for the
corresponding half-hour period within the same month. The final cleaned
dataset contained 99.844\% original readings, 0.012\% interpolated values,
and 0.145\% household-profile estimates.
% CHECK: the notebook also drops partial first/last months regardless of
% coverage (most of the ~12,000 months removed), and 13 households
% (5,561 -> 5,548) drop out entirely. Neither is stated here.

\subsection{Duplicate Removal}
\label{sec:duplicates}

Duplicate records were checked using the combination of household
identifier (LCLid) and timestamp (tstp). No duplicate household-time
records were found; therefore, no duplicate removal was necessary.

\subsection{Outlier Detection \& Treatment}
\label{sec:outliers}

Outliers were assessed using physical consumption limits and
household-specific usage patterns. No negative readings or values
exceeding the selected threshold of 11.5~kWh per half-hour were
identified. A household-level comparison using ten times the
99th-percentile consumption identified 34 unusual readings across 14
households. These were retained because they remained physically
plausible and could represent genuine high-consumption events.

\subsection{Data Type Correction}
\label{sec:types}

Electricity consumption values were converted from text to numerical
(DOUBLE) format, and the energy column was renamed kwh. Timestamp
fields were validated as temporal values, while household identifiers
were checked for consistency. Removing invalid records reduced the
initial dataset to 167,811,461 valid observations.

\subsection{Formatting \& Standardization}
\label{sec:formatting}

Timestamps were standardized using UTC/GMT conventions, ensuring
consistent 30-minute intervals without daylight-saving adjustments.
% CHECK: the notebook does not convert timestamps; it verifies they are
% already UTC/GMT (48 readings per household on every clock-change day).
% "Standardized" overstates what was done.
Inconsistent ACORN category labels (ACORN-U and ACORN-) were
standardized as Unknown. Incorrectly formatted bank holiday dates were
corrected, and column names were standardized to improve readability and
consistency across datasets.
% CHECK: the notebook also adds two missing December 2011 holidays
% (2011-12-26, 2011-12-27); six swapped dates were corrected.

\subsection{Categorical Data Encoding}
\label{sec:encoding}

Categorical variables were converted into numerical representations
suitable for machine learning. Tariff categories were encoded using
binary indicators, while ACORN socioeconomic classifications were
transformed using one-hot encoding. Original categorical values were
retained for interpretation.
% CHECK: only acorn_group (4 groups) is one-hot encoded; the detailed
% acorn category stays text.

\subsection{Feature Scaling / Normalization}
\label{sec:scaling}

Feature scaling was intentionally deferred to the modelling phase. This
allows scaling parameters to be estimated using only the appropriate
training data, preventing data leakage and preserving the original
measurement units during preprocessing.

\subsection{Data Integration}
\label{sec:integration}

Cleaned electricity readings were aggregated into daily and monthly
consumption records. Household characteristics, daily weather summaries,
and bank holiday indicators were integrated into the daily dataset.
Monthly electricity costs were estimated using a standard tariff of
\pounds 0.14228 per kWh. These estimates provide consistent billing targets but
do not represent actual invoices for households using time-of-use
tariffs.
% CHECK: the cost is a proxy for every household, not only ToU ones
% (no standing charge or VAT). As written, it implies Std households'
% figures are real invoices.

% =============================================================================
\section{Challenges and Solutions}
\label{sec:challenges}

The principal challenge was processing over 167 million records without
excessive memory consumption. DuckDB and Parquet addressed this through
efficient SQL-based processing and compressed columnar storage.
Additional challenges included invalid readings, missing intervals,
extended zero-consumption periods, inconsistent categorical labels, and
incorrect holiday dates. These were addressed through validation,
coverage-based filtering, interpolation, household-profile estimation,
and formatting corrections. Minor gaps in weather observations were
retained because their effect on daily averages was considered
negligible.

% =============================================================================
\section{Final Dataset and Validation}
\label{sec:final}

The preprocessing pipeline produced cleaned half-hourly, daily, and
monthly Parquet datasets. The cleaned half-hourly dataset contains
\textbf{155,510,640 records from 5,548 households}, covering 106,225
household-months. The daily dataset contains 3,239,805 household-day
records.

Appendix~\ref{app:outputs} lists every output file.

Final quality checks confirmed the absence of duplicate household-time
records, missing or negative energy values, and off-grid timestamps. The
resulting datasets provide a consistent foundation for electricity
consumption analysis and monthly bill forecasting.

\printbibliography[title={References}]

% =============================================================================
\appendix
\section{Output Files}
\label{app:outputs}

Files written by \code{notebooks/data\_preparation.ipynb}. Intermediate files are in
\code{data/interim/}; model-ready files are in \code{data/processed/}.

\begin{table}[h]
  \centering
  \caption{Output files.}
  \label{tab:outputs}
  \small
  % Y{f} widths must sum to 2 (the number of Y columns; Rows is fixed width).
  \begin{tabularx}{\linewidth}{Y{0.8} >{\raggedleft\arraybackslash}p{2cm} Y{1.2}}
    \tableheadrow
    \colhead{File} & \colhead{Rows} & \colhead{Contents} \\
    \tableheadend
    \multicolumn{3}{l}{\textit{\code{data/interim/}}} \\
    \code{halfhourly\_raw.parquet} & 167,817,021 & CSVs converted to Parquet unchanged (energy still text) \\
    \code{halfhourly\_typed.parquet} & 167,811,461 & \code{Null} rows dropped; energy cast to \code{kwh} (\code{DOUBLE}) \\
    \code{halfhourly\_grid.parquet} & 168,268,119 & Every expected half hour per household; missing slots have \code{kwh} = NULL \\
    \code{halfhourly\_clean.parquet} & 155,510,640 & Kept months only, all gaps filled; \code{status} marks each value as real, interpolated or profile-filled \\
    \multicolumn{3}{l}{\textit{\code{data/processed/}}} \\
    \code{daily.parquet} & 3,239,805 & Daily kWh and cost, real and filled slot counts, tariff, ACORN, temperature, holiday flag \\
    \code{monthly.parquet} & 106,225 & One row per household-month: kWh, bill, days, mean temperature, filled slots \\
    \code{monthly\_encoded.parquet} & 106,225 & Monthly table plus 0/1 columns for tariff and ACORN group \\
    \bottomrule
  \end{tabularx}
\end{table}

\end{document}
